% Thesis-only counterpart to desafios.tex (which frames open challenges as
% BitNets-vs-SNN columns -- not applicable once BitNet is out of the
% deck). Content drawn from the thesis' own "Limitações" section
% (documentation/00-thesis/monography/chapters/10-conclusions.tex), not
% invented for this talk: these are the three design limitations the
% thesis itself names, so the live software's demos (paraconsistent.*,
% comparison.autoencoders) are shown next to the caveats the real
% experiment ran into, not just the mechanics in isolation.
\begin{frame}
	\frametitle{Limitações do experimento}
	\small
	\begin{block}{A limitação dominante: tamanho da base}
		15 locutores, 1974 amostras. Sob protocolo de verificação --- partições disjuntas por locutor --- cada dobra treina com poucas identidades: risco alto de o classificador não ter sinal suficiente para generalizar.
	\end{block}
	\begin{itemize}
		\item As três categorias de representação (manual, \textit{autoencoder}, fusão) \textbf{não} foram comparadas por um classificador comum --- só pelo critério paraconsistente.
		\item Os \textit{autoencoders} tiveram $\bar\beta > 0{,}97$ em toda configuração: latentes perto do colapso. Não dá para distinguir, com os dados disponíveis, ``extrator pior'' de ``treino insuficiente''.
	\end{itemize}
	\begin{alertblock}{Consequência}
		A seleção paraconsistente ranqueia qualidade \textbf{antes} de treinar um classificador --- mas ranquear bem um extrator degenerado via treino insuficiente ainda é um risco que $D_{\text{penalizado}}$ por si só não resolve.
	\end{alertblock}
\end{frame}
